% Generated draft: requires mathematical and source review.
\documentclass[11pt]{article}
\usepackage[margin=25mm]{geometry}
\usepackage{fontspec}
\setmainfont{Latin Modern Roman}
\usepackage{amsmath,amssymb,mathtools,unicode-math}
\setmathfont{Latin Modern Math}
\usepackage{enumitem,needspace,xurl,hyperref}
\setlength{\parindent}{0pt}
\setlength{\parskip}{5pt}
\setlength{\emergencystretch}{3em}
\newcommand{\N}{\mathbb{N}}
\newcommand{\Z}{\mathbb{Z}}
\newcommand{\Q}{\mathbb{Q}}
\newcommand{\R}{\mathbb{R}}
\newcommand{\C}{\mathbb{C}}
\newcommand{\F}{\mathbb{F}}
\newcommand{\A}{\mathbb{A}}
\newcommand{\PP}{\mathbb{P}}
\newcommand{\E}{\mathbb{E}}
\newcommand{\eps}{\varepsilon}
\newcommand{\dd}{\,\mathrm{d}}
\begin{document}
\section*{UnsolvedMath ID: 30005200}
\textbf{Source problem code:} catalogue entry
\subsection{Definitions and mathematical statement}
A Kirchberg algebra means a nonzero, separable, simple, nuclear, purely infinite \(C^*\)-algebra. Let \(G\) be a locally compact group, as indicated by the supplied background, and let \(\alpha:G\curvearrowright A\) be an amenable action on a Kirchberg algebra \(A\). For a closed subgroup \(H\leq G\), write \(\alpha|_H\) for the restricted action and interpret the source notation \(A\rtimes H\) as \(A\rtimes_{\alpha|_H}H\).

The supplied question asks whether the following implication holds for every such action:
\[
\left[
\text{for every closed subgroup }H\leq G,
\quad A\rtimes_{\alpha|_H}H
\text{ is a Kirchberg algebra}
\right]
\quad\Longrightarrow\quad
\alpha\text{ is isometrically shift-absorbing}.
\]

\emph{Unresolved formulation qualifications.} The supplied record does not define amenability of the action or isometric shift-absorption, specify full versus reduced crossed products, or state continuity and countability hypotheses on the action and group. These conventions must be recovered from the cited report before this becomes a fully self-contained formulation. No second-countability assumption or additional restriction to discrete or compact groups is imposed here. The compact-group result mentioned in the supplied literature assessment is contextual information, not a restriction of the question.
\subsection{Short English statement}
For an amenable action of a locally compact group on a Kirchberg algebra, suppose restriction to every closed subgroup has a Kirchberg crossed product. Must the action be isometrically shift-absorbing? The precise conventions for the specialized properties remain to be confirmed from the source.
\subsection{Sources}
\begin{itemize}
\item Local UnsolvedMath registry, supplied record with permanent ID \(30005200\), title \emph{Characterization of Isometric Shift Absorption}, exact source problem code \texttt{catalogue entry}, in the Oberwolfach Reports Open Problems collection. The record gives proposed year 2022 and supplies no individual proposer.
\item Supplied report citation: \emph{\(C^*\)-Algebras (2023)}, Oberwolfach Reports, pp.~2059--2127, DOI \texttt{10.4171/owr/2022/36}. Original report: \url{https://doi.org/10.4171/owr/2022/36}. The supplied citation's year 2023 and its DOI year component 2022 are retained without reconciliation. The supplied original passage describes the broader aim of determining whether isometric shift-absorption is a familiar property in disguise, but does not itself include the hypotheses in the reconstructed statement.
\item Supplied literature reference: Mukohara--Suzuki, 2026, \url{https://arxiv.org/abs/2605.21128}. No paper title or precise theorem statement is supplied.
\item Literature assessment dated exactly \texttt{2026-08-22}: the supplied curation labels the problem partially solved and says that the general locally compact group question remains unresolved, while Mukohara and Suzuki confirm the corresponding isometric shift-absorption question for compact group actions. This is supplied dated curation, not an independently verified assessment; the exact hypotheses and scope of the cited result require verification.
\item The registry labels the statement \texttt{corrected\_verified} and reports that it was reconstructed and checked against the official Oberwolfach report. This is a supplied source-verification claim, not verification performed for this draft.
\end{itemize}
\end{document}
